\chapter{KHẢO SÁT VÀ PHÂN TÍCH YÊU CẦU}

\section{Khảo sát hiện trạng}

\subsection{Các công cụ dịch thuật hiện có}

Thị trường hiện có nhiều công cụ hỗ trợ dịch thuật, song chưa có giải pháp nào đáp ứng đầy
đủ nhu cầu doanh nghiệp sản xuất đa ngôn ngữ. Bảng~\ref{tab:survey} so sánh các công cụ phổ
biến nhất với yêu cầu của hệ thống đề xuất.

\begin{table}[H]
\centering
\caption{So sánh các công cụ dịch thuật hiện có}
\label{tab:survey}
\renewcommand{\arraystretch}{1.3}
\begin{tabularx}{\textwidth}{|l|X|X|X|X|X|}
\hline
\textbf{Tiêu chí}
  & \textbf{Google\newline Translate}
  & \textbf{DeepL}
  & \textbf{SDL\newline Trados}
  & \textbf{Matecat}
  & \textbf{Hệ thống\newline đề xuất} \\
\hline
Hỗ trợ DOCX          & Có  & Có  & Có  & Có  & \textbf{Có}  \\
\hline
Hỗ trợ XLSX          & Không & Không & Có  & Hạn chế & \textbf{Có}  \\
\hline
Hỗ trợ PPTX          & Không & Không & Có  & Có  & \textbf{Có}  \\
\hline
Hỗ trợ PDF scan (OCR)& Không & Không & Hạn chế & Không & \textbf{Có}  \\
\hline
Từ điển thuật ngữ    & Không & Có   & Có  & Có  & \textbf{Có (3 cấp)} \\
\hline
Đề xuất thuật ngữ\newline cộng đồng
                     & Không & Không & Không & Không & \textbf{Có}  \\
\hline
Hỗ trợ 10+ ngôn ngữ\newline châu Á
                     & Có   & Hạn chế & Có & Hạn chế & \textbf{Có}  \\
\hline
Giao diện web\newline nội bộ doanh nghiệp
                     & Không & Không & Không & Có & \textbf{Có}  \\
\hline
SSO doanh nghiệp     & Không & Không & Có  & Có  & \textbf{Có (Cognito)} \\
\hline
Chi phí              & Miễn phí/ API & Freemium & Cao  & Trả phí & Tự triển khai \\
\hline
\end{tabularx}
\end{table}

Từ kết quả khảo sát, có thể thấy rằng không công cụ nào đáp ứng đồng thời bốn yêu cầu cốt lõi:
(1)~hỗ trợ đầy đủ các định dạng tài liệu doanh nghiệp bao gồm PDF scan, (2)~hệ thống thuật ngữ
theo ngữ cảnh doanh nghiệp với quy trình duyệt cộng đồng, (3)~tích hợp SSO doanh nghiệp, và
(4)~giao diện web triển khai nội bộ. Đây là cơ sở để đồ án xây dựng hệ thống riêng.

\subsection{Quy trình dịch tài liệu thủ công tại doanh nghiệp}

Quy trình dịch tài liệu thông thường tại các doanh nghiệp đa ngôn ngữ vẫn diễn ra thủ công:
nhân viên sao chép nội dung từ tài liệu gốc vào Google~Translate hoặc DeepL, sau đó sao chép
lại kết quả vào tài liệu đích và chỉnh sửa thủ công. Quy trình này có các nhược điểm rõ ràng:

\begin{itemize}
  \item Tốn thời gian, đặc biệt với tài liệu nhiều trang.
  \item Mất định dạng (bảng, bảng tính, hình ảnh nhúng).
  \item Thiếu nhất quán thuật ngữ giữa các tài liệu và giữa các nhân viên.
  \item Không có lịch sử lưu trữ tài liệu đã dịch.
\end{itemize}

\section{Tổng quan chức năng}

\subsection{Biểu đồ use case tổng quát}

Hệ thống phục vụ ba nhóm tác nhân chính:

\begin{itemize}
  \item \textbf{Người dùng (User):} Nhân viên trong tổ chức sử dụng hệ thống để dịch tài liệu và quản
    lý thuật ngữ cá nhân.
  \item \textbf{Library Keeper:} Nhân viên phụ trách quản lý thư viện thuật ngữ chung của công
    ty, có quyền duyệt/từ chối đề xuất thuật ngữ.
  \item \textbf{Admin:} Quản trị viên hệ thống, có toàn quyền quản lý người dùng, cấu hình hệ
    thống và xem thống kê.
\end{itemize}

Agent:

\begin{figure}[H]
\centering
% Mermaid/PlantUML source cho use case tổng quan:
%
% @startuml
% left to right direction
% actor "Người dùng" as User
% actor "Library Keeper" as LK
% actor "Admin" as Admin
%
% rectangle "Hệ thống dịch tài liệu đa ngôn ngữ" {
%   usecase "Đăng nhập" as UC01
%   usecase "Dịch tài liệu" as UC02
%   usecase "Dịch văn bản" as UC03
%   usecase "Xem lịch sử dịch" as UC04
%   usecase "Quản lý thư viện cá nhân" as UC05
%   usecase "Đề xuất thuật ngữ" as UC06
%   usecase "Quản lý thư viện chung" as UC07
%   usecase "Duyệt đề xuất thuật ngữ" as UC08
%   usecase "Quản lý người dùng" as UC09
%   usecase "Xem thống kê từ khóa" as UC10
%   usecase "Chuyển đổi định dạng" as UC11
% }
%
% User --> UC01
% User --> UC02
% User --> UC03
% User --> UC04
% User --> UC05
% User --> UC06
% User --> UC11
%
% LK --> UC01
% LK --> UC07
% LK --> UC08
% LK --|> User
%
% Admin --> UC01
% Admin --> UC09
% Admin --> UC10
% Admin --|> LK
% @enduml
%
\fbox{\parbox{0.9\textwidth}{\centering\vspace{1cm}
\textit{[Hình: Biểu đồ use case tổng quát]}\\[0.5cm]
\textbf{Image Description:} Biểu đồ use case với ba tác nhân: Người dùng (User), Library Keeper
và Admin. Hệ thống có 11 use case chính: đăng nhập, dịch tài liệu, dịch văn bản, xem lịch sử
dịch, quản lý thư viện cá nhân, đề xuất thuật ngữ, quản lý thư viện chung, duyệt đề xuất, quản
lý người dùng, xem thống kê từ khóa, chuyển đổi định dạng. Admin kế thừa toàn bộ quyền của
Library Keeper; Library Keeper kế thừa toàn bộ quyền của User.
\vspace{0.5cm}}}
\caption{Biểu đồ use case tổng quát}
\label{fig:usecase_overview}
\end{figure}

\subsection{Biểu đồ use case phân rã: Dịch tài liệu}

Agent:

\begin{figure}[H]
\centering
\fbox{\parbox{0.9\textwidth}{\centering\vspace{1cm}
\textit{[Hình: Biểu đồ use case phân rã Dịch tài liệu]}\\[0.5cm]
\textbf{Image Description:} Phân rã use case ``Dịch tài liệu'' gồm các bước: tải tệp lên S3,
chọn ngôn ngữ nguồn (tùy chọn, có auto-detect), chọn một hoặc nhiều ngôn ngữ đích, chọn chế
độ thư viện (Không dùng / Thư viện chung / Thư viện cá nhân), kích hoạt dịch, tải kết quả
về. Có extension ``Kiểm tra PDF'' khi tệp là PDF: nếu là native-text thì ConvertAPI→DOCX;
nếu là scan thì OCR pipeline.
\vspace{0.5cm}}}
\caption{Biểu đồ use case phân rã -- Dịch tài liệu}
\label{fig:usecase_translate}
\end{figure}

\subsection{Biểu đồ use case phân rã: Quản lý thuật ngữ}

Agent:

\begin{figure}[H]
\centering
\fbox{\parbox{0.9\textwidth}{\centering\vspace{1cm}
\textit{[Hình: Biểu đồ use case phân rã Quản lý thuật ngữ]}\\[0.5cm]
\textbf{Image Description:} Phân rã use case ``Quản lý thuật ngữ'' gồm hai nhánh: (1)~Thư viện
cá nhân (User): thêm, sửa, xóa từ khóa cá nhân; đề xuất từ khóa vào thư viện chung.
(2)~Thư viện chung (Library Keeper/Admin): xem danh sách đề xuất đang chờ, duyệt hoặc từ chối
từng đề xuất, quản lý từ khóa đã được duyệt, cấu hình ngưỡng số người đề xuất tối thiểu.
\vspace{0.5cm}}}
\caption{Biểu đồ use case phân rã -- Quản lý thuật ngữ}
\label{fig:usecase_library}
\end{figure}

\subsection{Quy trình nghiệp vụ: Đề xuất thuật ngữ vào thư viện chung}

Agent:

\begin{figure}[H]
\centering
\fbox{\parbox{0.9\textwidth}{\centering\vspace{1cm}
\textit{[Hình: Quy trình đề xuất thuật ngữ]}\\[0.5cm]
\textbf{Image Description:} Quy trình nghiệp vụ dạng BPMN/Flowchart: Người dùng thêm từ khóa
vào Thư viện cá nhân → Người dùng nhấn ``Đề xuất lên thư viện chung'' → Hệ thống lưu vào
KeywordQueue → Hệ thống đếm số lượt đề xuất duy nhất; nếu chưa đủ ngưỡng thì chờ; nếu đủ
ngưỡng thì tạo KeywordSuggestion với status=pending và gửi thông báo cho Library Keeper →
Library Keeper duyệt hoặc từ chối → Nếu duyệt: status=approved, cập nhật GCS và Google
Translation glossary theo 44 cặp ngôn ngữ.
\vspace{0.5cm}}}
\caption{Quy trình nghiệp vụ đề xuất thuật ngữ vào thư viện chung}
\label{fig:bpmn_suggest}
\end{figure}

\section{Đặc tả chức năng}

\subsection{Đặc tả use case: UC-01 -- Dịch tài liệu}

\begin{table}[H]
\centering
\caption{Đặc tả use case UC-01: Dịch tài liệu}
\label{tab:uc01}
\renewcommand{\arraystretch}{1.3}
\begin{tabularx}{\textwidth}{|l|X|}
\hline
\textbf{Tên use case}     & UC-01: Dịch tài liệu \\
\hline
\textbf{Tác nhân}         & Người dùng đã đăng nhập \\
\hline
\textbf{Điều kiện trước}  & Người dùng đã đăng nhập hệ thống \\
\hline
\textbf{Điều kiện sau}    & Tệp dịch được lưu trên S3 và bản ghi TranslatedFile được tạo trong DB \\
\hline
\textbf{Mô tả}            & Người dùng tải tệp lên (DOCX/XLSX/PPTX/PDF/CSV/XLS), chọn ngôn ngữ và
                            chế độ thư viện, hệ thống dịch và trả về liên kết tải về \\
\hline
\textbf{Luồng chính}      &
\begin{enumerate}[leftmargin=*,nosep]
  \item Người dùng truy cập trang chủ.
  \item Người dùng tải tệp lên AWS~S3 qua widget.
  \item Người dùng chọn ngôn ngữ nguồn (hoặc để auto-detect).
  \item Người dùng chọn một hoặc nhiều ngôn ngữ đích.
  \item Người dùng chọn chế độ thư viện (Không dùng / Chung / Cá nhân).
  \item Người dùng nhấn ``Dịch''.
  \item Hệ thống tải tệp từ S3 về server tạm.
  \item Hệ thống thiết lập ngữ cảnh dịch (library\_mode + user\_id) qua thread-local.
  \item Hệ thống dịch từng ngôn ngữ đích, upload kết quả lên S3.
  \item Hệ thống lưu bản ghi TranslatedFile, trả về danh sách URL.
  \item Người dùng tải tệp về.
\end{enumerate} \\
\hline
\textbf{Luồng thay thế}  & 3a. Tệp là PDF: hệ thống gọi \texttt{is\_pdf\_truly\_editable()} để
                            phân loại, chuyển hướng sang pipeline phù hợp.
                            \newline
                            4a. Ngôn ngữ nguồn không xác định: hệ thống auto-detect từ nội dung
                            tệp. \\
\hline
\textbf{Ngoại lệ}        & ConvertAPI thất bại → trả lỗi 500; Upload S3 thất bại → trả lỗi 500 \\
\hline
\end{tabularx}
\end{table}

\subsection{Đặc tả use case: UC-02 -- Dịch văn bản}

\begin{table}[H]
\centering
\caption{Đặc tả use case UC-02: Dịch văn bản thuần}
\label{tab:uc02}
\renewcommand{\arraystretch}{1.3}
\begin{tabularx}{\textwidth}{|l|X|}
\hline
\textbf{Tên use case}    & UC-02: Dịch văn bản \\
\hline
\textbf{Tác nhân}        & Người dùng đã đăng nhập \\
\hline
\textbf{Điều kiện trước} & Người dùng đã đăng nhập \\
\hline
\textbf{Điều kiện sau}   & Văn bản dịch được hiển thị trên giao diện \\
\hline
\textbf{Mô tả}           & Người dùng nhập văn bản, chọn ngôn ngữ nguồn và đích, hệ thống trả
                           về văn bản đã dịch kèm thông tin ngôn ngữ \\
\hline
\textbf{Luồng chính}     &
\begin{enumerate}[leftmargin=*,nosep]
  \item Người dùng truy cập trang ``Dịch văn bản''.
  \item Người dùng nhập văn bản nguồn.
  \item Người dùng chọn ngôn ngữ nguồn (hoặc auto-detect).
  \item Người dùng chọn ngôn ngữ đích và chế độ thư viện.
  \item Hệ thống gọi \texttt{translate\_text\_with\_glossary()} với glossary\_id tương ứng.
  \item Hệ thống trả về văn bản dịch, ngôn ngữ nguồn phát hiện được.
\end{enumerate} \\
\hline
\textbf{Ngoại lệ}       & Ngôn ngữ đích không được hỗ trợ → lỗi 400; Lỗi API → lỗi 500 \\
\hline
\end{tabularx}
\end{table}

\subsection{Đặc tả use case: UC-03 -- Quản lý thư viện cá nhân}

\begin{table}[H]
\centering
\caption{Đặc tả use case UC-03: Quản lý thư viện cá nhân}
\label{tab:uc03}
\renewcommand{\arraystretch}{1.3}
\begin{tabularx}{\textwidth}{|l|X|}
\hline
\textbf{Tên use case}    & UC-03: Quản lý thư viện cá nhân \\
\hline
\textbf{Tác nhân}        & Người dùng \\
\hline
\textbf{Điều kiện trước} & Người dùng đã đăng nhập \\
\hline
\textbf{Điều kiện sau}   & Bản ghi PrivateKeyword được tạo/cập nhật/xóa trong DB; nếu cập nhật
                           từ có trong đề xuất thì GCS được đồng bộ bất đồng bộ \\
\hline
\textbf{Mô tả}           & Người dùng thêm, sửa hoặc xóa các cặp từ khóa cá nhân (PrivateKeyword)
                           theo 10 ngôn ngữ, có thể thêm ghi chú \\
\hline
\textbf{Luồng chính}     & Người dùng vào trang ``Thư viện cá nhân'' → thêm hàng mới hoặc chỉnh
                           sửa hàng hiện có → nhập giá trị từng ngôn ngữ → lưu.
                           Sau khi lưu, nếu người dùng có đề xuất đang chờ liên quan đến từ
                           khóa này, hệ thống trigger cập nhật glossary GCS qua background thread. \\
\hline
\end{tabularx}
\end{table}

\subsection{Đặc tả use case: UC-04 -- Quản lý thư viện chung}

\begin{table}[H]
\centering
\caption{Đặc tả use case UC-04: Quản lý thư viện chung (Library Keeper / Admin)}
\label{tab:uc04}
\renewcommand{\arraystretch}{1.3}
\begin{tabularx}{\textwidth}{|l|X|}
\hline
\textbf{Tên use case}    & UC-04: Quản lý thư viện chung \\
\hline
\textbf{Tác nhân}        & Library Keeper, Admin \\
\hline
\textbf{Điều kiện trước} & Người dùng có vai trò Library Keeper hoặc Admin \\
\hline
\textbf{Điều kiện sau}   & KeywordSuggestion được cập nhật status; GCS và Google Translation
                           glossary được đồng bộ bất đồng bộ qua background thread \\
\hline
\textbf{Mô tả}           & Xem danh sách đề xuất chờ duyệt (status=pending), duyệt hoặc từ chối
                           từng đề xuất, quản lý từ khóa đã phê duyệt (CRUD), cấu hình ngưỡng
                           \texttt{min\_suggesters\_for\_queue} \\
\hline
\end{tabularx}
\end{table}

\section{Yêu cầu phi chức năng}

\begin{table}[H]
\centering
\caption{Yêu cầu phi chức năng}
\label{tab:nonfunc}
\renewcommand{\arraystretch}{1.3}
\begin{tabularx}{\textwidth}{|l|l|X|}
\hline
\textbf{Loại}       & \textbf{Tên}          & \textbf{Mô tả / Tiêu chuẩn} \\
\hline
Hiệu năng           & Thời gian phản hồi    & Dịch tài liệu DOCX 10~trang $\leq 30$~giây; dịch văn bản
                                              $\leq 5$~giây \\
\hline
Hiệu năng           & Xử lý song song       & Hỗ trợ nhiều ngôn ngữ đích trong một yêu cầu (sequential loop
                                              trên server) \\
\hline
Bảo mật             & Xác thực              & JWT với TTL 15~phút (access) / 7~ngày (refresh); AWS Cognito
                                              OIDC cho môi trường production \\
\hline
Bảo mật             & Phân quyền            & Ba vai trò (User / Library Keeper / Admin) với kiểm tra
                                              \texttt{IsAuthenticated} trên mọi endpoint \\
\hline
Bảo mật             & Lưu trữ tệp           & Tệp lưu trên AWS~S3 với tên UUID ngẫu nhiên; không expose
                                              tên tệp gốc qua URL \\
\hline
Tin cậy             & Xử lý lỗi            & Pipeline dịch bắt ngoại lệ theo từng ngôn ngữ đích; kết quả
                                              thành công vẫn trả về dù một ngôn ngữ thất bại \\
\hline
Khả năng sử dụng    & Giao diện             & SPA trực quan với Ant Design; thông báo toast khi thành công/
                                              lỗi; hiển thị trạng thái đang dịch \\
\hline
Khả năng bảo trì    & Container hóa         & Toàn bộ backend và frontend Docker hóa; cấu hình qua biến
                                              môi trường \texttt{.env} \\
\hline
Khả năng mở rộng    & Ngôn ngữ              & Thêm ngôn ngữ mới chỉ cần bổ sung cặp glossary và cập nhật
                                              hằng số \texttt{language\_pair} \\
\hline
\end{tabularx}
\end{table}
